\documentclass[10pt,a4paper]{amsart}
\usepackage{amsmath,amssymb,amsthm,graphicx}
\usepackage{verbatim}
\newenvironment{reference}{\comment}{\endcomment}
%\newenvironment{reference}{}{}
\newenvironment{slogan}{\comment}{\endcomment}
\newenvironment{history}{\comment}{\endcomment}

% For commutative diagrams we use Xy-pic
\usepackage[all]{xy}

% We use 2cell for 2-commutative diagrams.
\xyoption{2cell}
% Theorem environments.
%
\theoremstyle{plain}
\newtheorem{theorem}[subsection]{Theorem}
\newtheorem{proposition}[subsection]{Proposition}
\newtheorem{lemma}[subsection]{Lemma}

\theoremstyle{definition}
\newtheorem{definition}[subsection]{Definition}
\newtheorem{example}[subsection]{Example}
\newtheorem{exercise}[subsection]{Exercise}
\newtheorem{situation}[subsection]{Situation}

\theoremstyle{remark}
\newtheorem{remark}[subsection]{Remark}
\newtheorem{remarks}[subsection]{Remarks}

\numberwithin{equation}{subsection}

% Macros
%
\def\lim{\mathop{\mathrm{lim}}\nolimits}
\def\colim{\mathop{\mathrm{colim}}\nolimits}
\def\Spec{\mathop{\mathrm{Spec}}}
\def\Hom{\mathop{\mathrm{Hom}}\nolimits}
\def\Ext{\mathop{\mathrm{Ext}}\nolimits}
\def\SheafHom{\mathop{\mathcal{H}\!\mathit{om}}\nolimits}
\def\SheafExt{\mathop{\mathcal{E}\!\mathit{xt}}\nolimits}
\def\Sch{\mathit{Sch}}
\def\Mor{\mathop{\mathrm{Mor}}\nolimits}
\def\Ob{\mathop{\mathrm{Ob}}\nolimits}
\def\Sh{\mathop{\mathit{Sh}}\nolimits}
\def\NL{\mathop{N\!L}\nolimits}
\def\CH{\mathop{\mathrm{CH}}\nolimits}
\def\proetale{{pro\text{-}\acute{e}tale}}
\def\etale{{\acute{e}tale}}
\def\QCoh{\mathit{QCoh}}
\def\Ker{\mathop{\mathrm{Ker}}}
\def\Im{\mathop{\mathrm{Im}}}
\def\Coker{\mathop{\mathrm{Coker}}}
\def\Coim{\mathop{\mathrm{Coim}}}

% Boxtimes
%
\DeclareMathSymbol{\boxtimes}{\mathbin}{AMSa}{"02}

%
% Macros for moduli stacks/spaces
%
\def\QCohstack{\mathcal{QC}\!\mathit{oh}}
\def\Cohstack{\mathcal{C}\!\mathit{oh}}
\def\Spacesstack{\mathcal{S}\!\mathit{paces}}
\def\Quotfunctor{\mathrm{Quot}}
\def\Hilbfunctor{\mathrm{Hilb}}
\def\Curvesstack{\mathcal{C}\!\mathit{urves}}
\def\Polarizedstack{\mathcal{P}\!\mathit{olarized}}
\def\Complexesstack{\mathcal{C}\!\mathit{omplexes}}
% \Pic is the operator that assigns to X its picard group, usage \Pic(X)
% \Picardstack_{X/B} denotes the Picard stack of X over B
% \Picardfunctor_{X/B} denotes the Picard functor of X over B
\def\Pic{\mathop{\mathrm{Pic}}\nolimits}
\def\Picardstack{\mathcal{P}\!\mathit{ic}}
\def\Picardfunctor{\mathrm{Pic}}
\def\Deformationcategory{\mathcal{D}\!\mathit{ef}}
\usepackage{cleveref}
\crefname{originalreference}{original reference}{original references}
\Crefname{originalreference}{Original reference}{Original references}
\makeatletter
\newlabel{sites-modules-lemma-finite-presentation-quasi-compact-colimit}{{0GN0}{source}{}{sourcepointer-1}{}}
\newlabel{sites-modules-lemma-finite-presentation-quasi-compact-colimit@cref}{{[originalreference][0][]0GN0}{[1][0][]source}{}{}{}}
\makeatother
\crefname{originalreference}{original reference}{original references}
\Crefname{originalreference}{Original reference}{Original references}
\makeatletter
\newlabel{sites-lemma-directed-colimits-global-sections}{{0GMR}{source}{}{sourcepointer-2}{}}
\newlabel{sites-lemma-directed-colimits-global-sections@cref}{{[originalreference][0][]0GMR}{[1][0][]source}{}{}{}}
\makeatother
\crefname{originalreference}{original reference}{original references}
\Crefname{originalreference}{Original reference}{Original references}
\makeatletter
\newlabel{sites-remark-stronger-conditions}{{0GMS}{source}{}{sourcepointer-3}{}}
\newlabel{sites-remark-stronger-conditions@cref}{{[originalreference][0][]0GMS}{[1][0][]source}{}{}{}}
\makeatother
\crefname{originalreference}{original reference}{original references}
\Crefname{originalreference}{Original reference}{Original references}
\makeatletter
\newlabel{stacks-properties-lemma-quasi-compact-stack}{{04YC}{source}{}{sourcepointer-4}{}}
\newlabel{stacks-properties-lemma-quasi-compact-stack@cref}{{[originalreference][0][]04YC}{[1][0][]source}{}{}{}}
\makeatother
\crefname{originalreference}{original reference}{original references}
\Crefname{originalreference}{Original reference}{Original references}
\makeatletter
\newlabel{topologies-lemma-fppf-affine}{{021P}{source}{}{sourcepointer-5}{}}
\newlabel{topologies-lemma-fppf-affine@cref}{{[originalreference][0][]021P}{[1][0][]source}{}{}{}}
\makeatother
\crefname{originalreference}{original reference}{original references}
\Crefname{originalreference}{Original reference}{Original references}
\makeatletter
\newlabel{stacks-morphisms-lemma-quasi-compact-quasi-separated-permanence}{{075S}{source}{}{sourcepointer-6}{}}
\newlabel{stacks-morphisms-lemma-quasi-compact-quasi-separated-permanence@cref}{{[originalreference][0][]075S}{[1][0][]source}{}{}{}}
\makeatother
\crefname{originalreference}{original reference}{original references}
\Crefname{originalreference}{Original reference}{Original references}
\makeatletter
\newlabel{stacks-sheaves-lemma-quasi-coherent-algebraic-stack}{{06WV}{source}{}{sourcepointer-7}{}}
\newlabel{stacks-sheaves-lemma-quasi-coherent-algebraic-stack@cref}{{[originalreference][0][]06WV}{[1][0][]source}{}{}{}}
\makeatother
\title{Hom out of a finitely presented module commutes with filtered colimits on a qcqs stack}
\author{The Stacks Project Authors}
\date{stacks-project:stacks-cohomology.tex:local-snapshot-20261012; GNU FDL 1.2 or later}
\begin{document}
\maketitle
\noindent\textit{Editorial note.} All delivered mathematical text is the author's own, copied verbatim from the pinned Stacks Project chapter file stacks-cohomology.tex; no statement or proof step was rewritten, reordered or removed.\par\medskip
\noindent\textit{Editorial note.} The counted span cites the lemmas recorded in assembly.external\_references with canonical tags as navigation pointers.\par\medskip
\noindent\textit{Allowed standard prerequisites (editorial).}\begin{itemize}
\item Finitely presented modules and filtered colimits. Modules on Sites, Lemma sites-modules-lemma-finite-presentation-quasi-compact-colimit. Module-level statement being instantiated.
\item Directed colimits of global sections, hypotheses. Sites, Lemma sites-lemma-directed-colimits-global-sections (4) with Remark sites-remark-stronger-conditions. Provides the conditions (2)(a)-(c) verified.
\item QCoh on an algebraic stack. Sheaves on Stacks, Lemma stacks-sheaves-lemma-quasi-coherent-algebraic-stack. Colimit commutation of the inclusion; fp modules are qc.
\end{itemize}\medskip
\ifdefined\hypertarget\hypertarget{sourcepointer-1}{}\fi
\noindent\textit{Source reference.} Original reference 0GN0: sources/stacks-chow/tags.txt, source lines 20775--20775. Cross-chapter standard prerequisite cited by the author inside the counted span; it is used as a pointer and its own proof is not part of this item.\par\medskip
\ifdefined\hypertarget\hypertarget{sourcepointer-2}{}\fi
\noindent\textit{Source reference.} Original reference 0GMR: sources/stacks-chow/tags.txt, source lines 20766--20766. Cross-chapter standard prerequisite cited by the author inside the counted span; it is used as a pointer and its own proof is not part of this item.\par\medskip
\ifdefined\hypertarget\hypertarget{sourcepointer-3}{}\fi
\noindent\textit{Source reference.} Original reference 0GMS: sources/stacks-chow/tags.txt, source lines 20767--20767. Cross-chapter standard prerequisite cited by the author inside the counted span; it is used as a pointer and its own proof is not part of this item.\par\medskip
\ifdefined\hypertarget\hypertarget{sourcepointer-4}{}\fi
\noindent\textit{Source reference.} Original reference 04YC: sources/stacks-chow/tags.txt, source lines 6425--6425. Cross-chapter standard prerequisite cited by the author inside the counted span; it is used as a pointer and its own proof is not part of this item.\par\medskip
\ifdefined\hypertarget\hypertarget{sourcepointer-5}{}\fi
\noindent\textit{Source reference.} Original reference 021P: sources/stacks-chow/tags.txt, source lines 2678--2678. Cross-chapter standard prerequisite cited by the author inside the counted span; it is used as a pointer and its own proof is not part of this item.\par\medskip
\ifdefined\hypertarget\hypertarget{sourcepointer-6}{}\fi
\noindent\textit{Source reference.} Original reference 075S: sources/stacks-chow/tags.txt, source lines 9141--9141. Cross-chapter standard prerequisite cited by the author inside the counted span; it is used as a pointer and its own proof is not part of this item.\par\medskip
\ifdefined\hypertarget\hypertarget{sourcepointer-7}{}\fi
\noindent\textit{Source reference.} Original reference 06WV: sources/stacks-chow/tags.txt, source lines 8829--8829. Cross-chapter standard prerequisite cited by the author inside the counted span; it is used as a pointer and its own proof is not part of this item.\par\medskip
\begin{lemma}
\label{lemma-finite-presentation-quasi-compact-colimit}
Let $\mathcal{X}$ be a quasi-compact and quasi-separated algebraic stack.
Let $I$ be a directed set and let $(\mathcal{F}_i, \varphi_{ii'})$ be a
system over $I$ of $\mathcal{O}_\mathcal{X}$-modules. Let $\mathcal{G}$ be an
$\mathcal{O}_\mathcal{X}$-module of finite presentation. Then we have
$$
\colim_i \Hom_\mathcal{X}(\mathcal{G}, \mathcal{F}_i)
=
\Hom_\mathcal{X}(\mathcal{G}, \colim_i \mathcal{F}_i).
$$
In particular, $\Hom_\mathcal{X}(\mathcal{G}, -)$ commutes with filtered
colimits in $\QCoh(\mathcal{O}_\mathcal{X})$.
\end{lemma}

\begin{proof}
The displayed equality is a special case of Modules on Sites, Lemma
\ref{sites-modules-lemma-finite-presentation-quasi-compact-colimit}.
In order to apply it, we need to check the hypotheses of
Sites, Lemma \ref{sites-lemma-directed-colimits-global-sections} part (4)
for the site $\mathcal{X}_{fppf}$. In order to do this, we will check
hypotheses (2)(a), (2)(b), (2)(c) of
Sites, Remark \ref{sites-remark-stronger-conditions}.
Namely, let $\mathcal{B} \subset \Ob(\mathcal{X}_{fppf})$ be the set
of objects lying over affine schemes. In other words, an element
of $\mathcal{B}$ is a morphism $x : U \to \mathcal{X}$ with $U$ affine.
We check each of the conditions (2)(a), (2)(b), and (2)(c)
of the remark in turn:
\begin{enumerate}
\item Since $\mathcal{X}$ is quasi-compact, there exists a
surjective and smooth morphism $x : U \to \mathcal{X}$ with $U$ affine
(Properties of Stacks, Lemma
\ref{stacks-properties-lemma-quasi-compact-stack}).
Then $h_x^\# \to *$ is a surjective map of sheaves on $\mathcal{X}_{fppf}$.
\item Since coverings in $\mathcal{X}_{fppf}$ are fppf
coverings, we see that every covering of $U \in \mathcal{B}$ is refined
by a finite affine fppf covering, see
Topologies, Lemma \ref{topologies-lemma-fppf-affine}.
\item Let $x : U \to \mathcal{X}$ and $x' : U' \to \mathcal{X}$ be
in $\mathcal{B}$. The product $h_x^\# \times h_{x'}^\#$ in
$\Sh(\mathcal{X}_{fppf})$ is equal to the sheaf on $\mathcal{X}_{fppf}$
determined by the algebraic space $W = U \times_{x, \mathcal{X}, x'} U'$ over
$\mathcal{X}$: for an object $y : V \to \mathcal{X}$ of
$\mathcal{X}_{fppf}$ we have
$(h_x^\# \times h_{x'}^\#)(y) = \{f : V \to W \mid y =
x \circ \text{pr}_1 \circ f = x' \circ \text{pr}_2 \circ f\}$.
The algebraic space $W$ is quasi-compact because $\mathcal{X}$
is quasi-separated, see Morphisms of Stacks, Lemma
\ref{stacks-morphisms-lemma-quasi-compact-quasi-separated-permanence}
for example. Hence we can choose an affine scheme $U''$ and
a surjective \'etale morphism $U'' \to W$. Denote $x'' : U'' \to \mathcal{X}$
the composition of $U'' \to W$ and $W \to \mathcal{X}$. Then
$h_{x''}^\# \to h_x^\# \times h_{x'}^\#$ is surjective as desired.
\end{enumerate}
For the final statement, observe that the inclusion functor
$\QCoh(\mathcal{O}_X) \to \textit{Mod}(\mathcal{O}_X)$
commutes with colimits and that finitely presented modules
are quasi-coherent. See Sheaves on Stacks, Lemma
\ref{stacks-sheaves-lemma-quasi-coherent-algebraic-stack}.
\end{proof}

\end{document}
